\documentclass[twoside]{article}
\usepackage{graphicx}
\usepackage{polski}
\usepackage{fancyhdr}
\usepackage{lastpage}
\usepackage{epigraph}
\usepackage{listings}
\usepackage{soul}
\usepackage{color}
\usepackage{multirow}
\usepackage{subcaption}
\usepackage{fancyvrb}
\usepackage{hyperref}


\title{AiSD Ex. 6 (Pon.)}
\author{\hl{Mateusz Kamieniecki}}
\date{\today}

\lstset{
    numbers=left
}

\makeatletter

\begin{document}

\pagestyle{fancy}
\fancyhead{}
\fancyhead[L]{\@title}
\fancyhead[R]{\@author}
\fancyhead[C]{\@date}

\cfoot{\thepage\ / \pageref{LastPage}}


\newpage

\section{Wprowadzenie}
Niektóre problemy wymagają od nas znalezienia rozwiązania pośród wielu możliwych rozwiązań -- są
to problemy optymalizacyjne. Szukając najlepszego rozwiązania (rozwiązania optymalnego), chcemy
znaleźć takie rozwiązanie, w którym końcowy wynik (w zależności od zagadnienia) będzie największy
lub najmniejszy. Co istotne, rozwiązanie danego problemu może mieć więcej niż jedno rozwiązanie
optymalne (dające ten sam wynik).

Na podstawie poprzednich zajęć wiemy, że niektóre problemy optymalizacyjne można rozwiązywać
wykorzystując podejście zachłanne (najpierw wybieramy podproblem, który w danym momencie wydaje się najlepszy, a następnie go rozwiązujemy), w którym rozwiązywany problem musi mieć własności:
optymalnej podstruktury i wyboru zachłannego.

Innym -- dość podobnym choć ogólniejszym -- podejściem do problemów optymalizacyjnych jest 
programowanie dynamiczne.
Technika ta zakłada (podobnie jak w podejściu zachłannym), że nasze optymalne rozwiązania problemów
składają się z optymalnych rozwiązań podproblemów, a same podproblemy są od siebie niezależne będąc
jednocześnie wspólnymi podproblemami dla innych problemów.

Jak to rozumieć?
Poszukajmy porównania do podejścia \textit{dziel i zwyciężaj}.
Dla ustalenia uwagi przyjmijmy, że rozpatrujemy algorytm szybkiego sortowania (algorytm 
typu \textit{dziel i zwyciężaj}). 
W tym zagadnieniu mieliśmy duży problem (posortowanie całej tablicy -- ustalmy, że sortowanie 
dotyczyło tablicy liczb całkowitych), który dzieliliśmy na mniejsze podproblemy (mniejsze podproblemy dzieliliśmy na jeszcze mniejsze podproblemy, i tak postępowaliśmy, aż osiągnęliśmy
ustalony rozmiar podproblemu).
Te podproblemy były od siebie niezależne.
Rozwiązanie (posortowanie) części tablicy nie wpływało na wynik posortowania innej części tablicy.
Jednocześnie mając podproblem \verb|pp1|, nie mogliśmy w żaden sposób wykorzystać
rozwiązania jego podproblemów (\verb|pp1_1|, \verb|pp1_2|) jako rozwiązań innych podproblemów
(np.: \verb|pp2|).

W przypadku programowania dynamicznego, mimo że problemy (podproblemy) są niezależne, 
to jednocześnie są one wspólne, czyli na drodze do znalezienia rozwiązania głównego problemu,
wielokrotnie rozwiązujemy te same (mniejsze) podproblemy.

A dlaczego to jest takie istotne?
Wynika to ze sposobu konstrukcji rozwiązań w programowaniu dynamicznym: zaczynamy rozwiązywać
duże problemy od ich małych podproblemów, których rozwiązania są zapamiętywane.
Jeśli jakiś podproblem wystąpi ponownie, to zamiast wykonywać kolejny raz te same obliczenia, 
odwołujemy się do zapamiętanych wcześniej rozwiązań.
Co oznacza, że korzystna jest sytuacja, gdy wielokrotnie w rozwiązaniu głównego problemu występują
te same podproblemy.
W takim zagadnieniu (naiwnym, bez techniki programowania dynamicznego), podejście
rekurencyjne może znacząco wydłużyć czas poszukiwania prawidłowego rozwiązania.

Spróbujmy to zobaczyć na poniższym przykładzie (inspiracja: C. L. R. S.).


\subsection{Problem cięcia pręta}

Problem cięcia pręta wygląda następująco: fabryka produkuje pręt o długości \textit{L}, a następnie
chce na nim jak najwięcej zarobić.
Okazuje się, że fragmenty pręta o różnej długości mają różną cenę.
W zależności od tego, na jakie fragmenty potniemy oryginalny pręt, fabryka zarobi na tym jednym
pręcie mniej lub więcej pieniędzy.

\begin{table}[h!]
\centering
\begin{tabular}{|l|r|r|r|r|r|r|r|r|r|r|}
\hline
Długość pręta $L$ [m] & 1 & 2 & 3 & 4  & 5  & 6  & 7  & 8  & 9  & 10 \\ \hline
Cena pręta $c_L$ [zł] & 2 & 6 & 9 & 10 & 11 & 18 & 18 & 21 & 25 & 31 \\ \hline
\end{tabular}
\caption{Tabela pokazująca cenę (zł) pręta o długości $L$ metrów.}
\label{table:rod_price}
\end{table}

Jeśli przyjmiemy, że znana jest cena $c_L$ pręta o długości $L$ -- co pokazano w~tabeli 
\ref{table:rod_price}. -- to możemy ustalić wszystkie rozwiązania (wszystkie możliwe zestawy cięć),
z których wybierzemy najlepsze (dające największy zysk).
Przykładowo: jeśli pręt ma długość 10 metrów, a cięcie jest możliwe tylko na odcinki pełnych metrów,
to możemy wykonać cięcie na $2^{L-1}$ sposobów (choć część rozwiązań będzie powtórzonych, ponieważ cięcie pręta od lewej strony jest analogicznym cięciem pręta od prawej strony).

Problem cięcia pręta jest problemem charakteryzującym się optymalną podstrukturą:
optymalne rozwiązanie problemu (pocięcia pręta) wymaga znalezienia optymalnego rozwiązania 
podproblemów (optymalnego pocięcia mniejszych fragmentów oryginalnego pręta) dających jak największy
zysk dla dwóch pociętych fragmentów.

\begin{table}[h]
\centering
\begin{tabular}{|r|r|r|l|}
\hline
Lp. & Długość pręta & Optymalne rozwiązanie & Cięcia                     \\ \hline
1.  & 1             & 2                     & -                          \\ \hline
2.  & 2             & 6                     & -                          \\ \hline
3.  & 3             & 9                     & -                          \\ \hline
4.  & 4             & 12                    & 2+2                        \\ \hline
5.  & 5             & 15                    & 2+3                        \\ \hline
6.  & 6             & 18                    & bez cięć; 3+3; 2+2+2       \\ \hline
7.  & 7             & 21                    & 2+2+3                      \\ \hline
8.  & 8             & 24                    & 2+6; 2+3+3                 \\ \hline
9.  & 9             & 27                    & 3+3+3; 3+6; 2+2+2+3        \\ \hline
10. & 10            & 31                    & -                          \\ \hline
\end{tabular}
\caption{Zestawienie optymalnych wyników dla problemu cięcia pręta o długości od 1 do 10.}
\end{table}

Rekurencyjnie można ten problem przedstawić opisowo w następujący sposób:
wybieramy pierwszą część, a dla pozostałej części próbujemy odszukać (rekurencyjnie) optymalne 
rozwiązanie.
Jeśli przez $r$ oznaczymy oczekiwane maksymalne rozwiązanie, to optymalne rozwiązanie dla 
pręta o długości $L$ można zapisać jako:
$$r_L=max(c_i+r_{L-i})$$, dla $i \leq L$.

Problem w tym, że takie podejście, choć końcowo powinno przynieść oczekiwany wynik, jest kosztowne
obliczeniowo -- zapisanie tego wzoru w najprostszej rekurencyjnej implementacji będzie 
charakteryzowało się złożonością wykładniczą $O(2^L)$.

Dlatego też ten problem sensowniej jest rozwiązać wykorzystując programowanie dynamiczne.
Do wyboru są dwa podejścia:
\begin{itemize}
    \item metoda zstępująca ze spamiętywaniem, czyli podejście rekurencyjne wykorzystujące 
    zapamiętane wcześniejsze wyniki obliczeń);
    \item metoda wstępująca -- zaczynamy od najmniejszych podproblemów i stopniowo rozwiązujemy 
    coraz większe (oczekujące wyników z mniejszych, już obliczonych podproblemów).
\end{itemize}

Dla obu podejść do problemu cięcia pręta złożoność czasowa powinna być klasy $O(L^2)$ 
(a nawet $\theta(L^2)$).

Można zauważyć, że w ogólnym problemie programowania dynamicznego, jeśli wszystkie możliwe 
podproblemy trzeba rozwiązać przynajmniej jeden raz, to metoda wstępująca powinna dawać lepszy wynik
(unikamy obsługi wywołań rekurencyjnych).
W przypadku odwrotnym, gdy nie wszystkie możliwe podproblemy wymagają rozwiązania (niektóre 
istniejące podproblemy są pomijane, czyli ich wyznaczanie nie przekłada się na rozwiązanie 
większych problemów), to metoda zstępująca (podejście rekurencyjne ze spamiętywaniem) może okazać
się efektywniejsza, ponieważ będzie zajmowała się tylko tymi podproblemami, które rzeczywiście są
konieczne do znalezienia rozwiązania optymalnego (a pozostałe pominie).


\subsection{Problem nawiasowania iloczynu macierzy}

Problem mnożenia ciągu macierzy jest problemem nawiasowania, czyli ustalenia kolejności mnożeń
poszczególnych macierzy w taki sposób, aby zminimalizować liczbę koniecznych mnożeń skalarnych.
Zasadniczo w tym problemie nie chodzi o wykonanie samych obliczeń (wymnożenie macierzy), a 
,,jedynie'' o ustalenie kolejności, w jakiej powinny być ze sobą mnożone.
Optymalne nawiasowanie iloczynu $N$ ($N >= 2$) macierzy wymaga znalezienia optymalnego nawiasowania 
,,podiloczynów'' (podproblemów).

Powinniśmy zauważyć, że koszt wymnożenia macierzy $M_i, M_{i+1}, ..., M_j$ (dla $i < j$) wymaga 
znalezienia optymalnego iloczynu podproblmów: $(M_i, M_{i+1}, ..., M_k)$ oraz 
$(M_{k+1}, M_{k+2}, M_j)$ (dla $i \leq k < j$), 
a następnie konieczne jest wymnożenie wyniku tych podproblemów ze sobą.

Czy to rzeczywiście ma sens?
Wiemy, że musimy wymnożyć ze sobą ciąg macierzy. Czyli musimy pomnożyć ze sobą dwie macierze (które 
mogły powstać w wyniku mnożenia innych macierzy). Szukamy takiej kombinacji, która da nam optymalny 
(w tym przypadku najmniejszy) wynik liczby mnożeń. Znajdując docelowe (optymalne) macierze (czyli 
rozwiązania podproblemów), wiemy już, że one (te podproblemy) też musiały być optymalnymi rozwiązania dla swoich podproblemów.
Nie może być inaczej, prawda?

Zatem docieramy do znanej nam już własności optymalnych podproblemów (w programowaniu dynamicznym).
Powinniśmy oczywiście pamiętać, że takich optymalnych rozwiązań może być więcej niż jedno, ale nas w pełni usatysfakcjonuje jedno z nich.

Okazuje się, że konieczność przeszukania wszystkich możliwych kombinacji nawiasowania iloczynu 
macierzy ma złożoność obliczeniową rzędu $\Omega(2^N)$, co oczywiście nie jest dla nas korzystnym 
sposobem rozwiązania problemu.
Implementując ten problem w wersji wstępującej można osiągnąć złożoność klasy $O(N^3)$. 
\newpage

\section{Zadania}

\subsection{Zadanie 0}
Proszę, analogicznie jak przy poprzednich ćwiczeniach,
uporządkować swoje repozytorium, aby organizacja projektów wyglądała jak poniżej. 
Oczywiście w tej strukturze mogą znajdować się również inne katalogi i pliki, ale proszę
dokładnie odwzorować przedstawiony porządek projektów -- dla uproszczenia zapisu w tym dokumencie
wyświetlanie poprzednich ćwiczeń zostało zredukowane do wyświetlania głównych katalogów.

\begin{verbatim}
.korzeń_repozytorium
+-- CW_0
+-- CW_1
+-- CW_2
+-- CW_3
+-- CW_4
+-- CW_5
+-- CW_6
    +-- AiSD2026ZEx6
        +-- pom.xml        
\end{verbatim}

\newpage

\subsection{Zadanie 1}
Dla dwóch problemów (cięcia pręta, mnożenia łańcucha macierzy) proszę wyznaczyć wielkości 
problemów (tabele: \ref{table:rod_cutter_results}, \ref{table:matrix_chain_results},
które mogą być realizowane przez poszczególne implementacje.

Jak ustalić wartość graniczną? Jeżeli czas potrzebny do rozwiązania problemu o wielkości N (czyli 
czas cięcia pręta o długości N lub czas mnożenia łańcucha N macierzy) zbliża się w przybliżeniu
do 4 sekund, to można uznać, że znalezione N jest naszą wartością graniczną.
Ten czas (4 sekundy) można dowolnie wydłużyć, ale na potrzeby naszego badania będzie 
to wartość w zupełności wystarczająca.

\begin{table}[h]
\centering
\begin{tabular}{|l|p{4cm}|}
\hline
\textbf{Implementacja}         & \textbf{Wielkość problemu, czyli długość pręta (N)} \\ \hline
Rekurencyjna                   &                                                     \\ \hline
Zstępująca (ze spamiętywaniem) &                                                     \\ \hline
Wstępująca                     &                                                     \\ \hline
\end{tabular}
\caption{Tabela pokazuje wielkości problemów (cięcia pręta), które można rozwiązać w sensownym czasie danym podejściem (rekurencyjnym, zstępującym, wstępującym).}
\label{table:rod_cutter_results}
\end{table}

\begin{table}[h]
\centering
\begin{tabular}{|l|p{4cm}|}
\hline
\textbf{Implementacja}         & \textbf{Wielkość problemu, czyli liczba macierzy (N)} \\ \hline
Rekurencyjna                   &                                                       \\ \hline
Zstępująca (ze spamiętywaniem) &                                                       \\ \hline
Wstępująca                     &                                                       \\ \hline
\end{tabular}
\caption{Tabela pokazuje wielkości problemów (nawiasowania mnożenia macierzy), które można rozwiązać w sensownym czasie danym podejściem (rekurencyjnym, zstępującym, wstępującym).}
\label{table:matrix_chain_results}
\end{table}

\textbf{Komentarz (\underline{przemyślane} wnioski):} \hl{TODO..................}

\subsection{Zadanie 2}
Proszę uzupełnić implementację metody \verb|reconstructOptimalSolutions()| z~klasy 
\verb|MatrixChainOrderExtendedResult|. 
Metoda powinna zwracać przykładowe (prawidłowe) odtworzenie rozwiązania (czyli tekstową 
reprezentację przedstawiającą kolejność mnożenia macierzy w łańcuchu).
Przykładowe uruchomienie (wynik) wraz z zaimplementowanym kodem źródłowym tej metody należy
umieścić w sprawozdaniu.

\subsection{Zadanie 3}
Proszę wykonać implementację (w klasie 
\verb|pl.edu.pw.ee.aisd2026zex6.lcs.| \verb|LongestCommonSubsequence|) 
problemu znajdowania najdłuższego wspólnego 
podciągu zgodnie z ustaleniami podanymi ,,na tablicy'' -- jeśli jeszcze nie były podane, to proszę 
się po nie zgłosić (w trakcie zajęć).

\section{Na zakończenie}
To już prawie koniec, ale ciągle nie wolno nam zapomnieć, że:
\begin{itemize}
    \item pod koniec zajęć zapisujemy swoje zmiany (robimy \textit{commit}) i oznaczamy tagiem ze znacznikiem czasowym \textbf{CW\_6\_EOT};
    \item wszystkie zmiany (cały kod źródłowy) muszą trafić do repozytorium.    
    A~ostatnią zmianę (\textit{commit}) należy oznaczyć tagiem ze znacznikiem czasowym \textbf{CW\_6\_FINAL};    
    \item końcową wersję tego dokumentu (PDF) należy umiecić w ISOD-zie.
\end{itemize}



\end{document}
